\section{Configurable Systems}
We are given a configuration system, and we need to reason about the configurations that contain the least amount of selected features.


\begingroup
\newcommand\fn[2]{\mathtt{#1}(#2)}
\newcommand\hi[1]{\fn{dec}{#1, 1}}
\newcommand\lo[1]{\fn{dec}{#1, 0}}



\subsection{Solution}
We identify the following parameters of the problem:
\begin{itemize}
  \item $n \in \mathbb{N}$ is the number of features, given in the header of the DIMACS file.
  \item $\varphi$ is the CNF formula from the DIMACS file. It consists of clauses $C_1,\dots,C_m$, and every clause is a set of literals over the features.
  \item $S=\{x_{218},x_{311}\}$ is the set of essential features.
\end{itemize}
We introduce the following variables:
\begin{itemize}
  \item $x_i \in \mathbb{B}$ for $i=1,\dots,n$: the value of $x_i$ is true if and only if feature $i$ is selected. The BDD variable order is given in the DIMACS file
  \item $\mathit{var}(u)$ and $\mathit{level}(u)\in\{1,\dots,n+1\}$ are the variable and the level of a BDD node $u$, so $\mathit{level}(x_i)=i$. The terminals $\top$ and $\bot$ are at level $n+1$. The children of $u$ are $\hi{u}$ and $\lo{u}$.
  \item $\mathtt{min\_cost}(u)\in\mathbb{N}\cup\{\infty\}$ is the smallest number of true variables on a path from $u$ to $\top$.
\end{itemize}

\subsubsection{Building the BDD}
We are given a DIMACS file, which contains an CNF.
We parse the CNF and after that we can build the BDD by iterating over the clauses, combining the variables with $\vee$ and combining the clauses with $\wedge$. 

\subsubsection{Computing the minimum amount of variables needed}
To compute the minimum amount of work we can 'walk' through the BDD, working from an bottoms-up manner we count the amount of 'true' edges we see.
This can be formulised in the function
\[
\fn{min\_cost}{n} =
\begin{cases}
\infty, & n = \bot \\
0, & n = \top \\
\min\bigl( \fn{min\_cost}{\lo{n}}, 1 + \fn{min\_cost}{\hi{n}} \bigr) & otherwise
\end{cases}
\]


\subsubsection{Computing the minimum amount of variables needed with essential variables}
Certain variables can be deemed as essentials. Those variables \textit{must} be true.
We achieve that though using the already build BDD and asserting those essential formulas.
The essential variables are: $\{x_{218}, x_{311}\}$. This produces the BDD:
\[
F \wedge x_{218} \wedge x_{311}
\]
Where $F$ is the BDD given in the DIMACS file.

\subsubsection{Building the BDD for the function with the minimum amount}
We again 'walk' through the BDD, but instead of computing a simple statistics we now produce a new BDD.
We do this via the \verb|min_config| function
\[
\fn{min\_config}{n} =
\begin{cases}
\bot & n = \bot \\
\top & n = \top \\
(\fn{var}{n} \wedge \fn{min\_hi}{n,\hi{n}}) \vee (\neg \fn{var}{n} \wedge \fn{min\_lo}{n,\lo{n}}) & otherwise
\end{cases}
\]
\[
\fn{min\_hi}{n, m} =
\begin{cases}
\bot & \fn{min\_cost}{m} \ge \fn{min\_cost}{n} \\
\fn{min\_config}{m} \wedge \fn{skip}{n,m}  & otherwise
\end{cases}
\]
\[
\fn{min\_lo}{n, m} =
\begin{cases}
\bot & \fn{min\_cost}{m} > \fn{min\_cost}{n} \\
\fn{min\_config}{m} \wedge \fn{skip}{n,m}  & otherwise
\end{cases}
\]
\[
  \fn{skip}{u,w} =\bigwedge_{\mathit{level}(u) < i < \mathit{level}(w)} \neg x_i 
\]

After we have restricted the BDD to only the configurations with the minimum amount, we also need to constrict the variables that may have been removed from the BDD.
We do this by collecting all the used variables in the BDD and asserting them all to false.
\[
\fn{used\_vars}{n} =
\begin{cases}
\emptyset & n = \bot \vee n = \top \\
\{\fn{var}{n}\} \cup \fn{used\_vars}{\hi{n}} \cup \fn{used\_vars}{\lo{n}} & otherwise
\end{cases}
\]
\[
\fn{constrict\_to\_used\_vars}{n,vars} = n \bigwedge_{v \in vars\backslash \fn{used\_vars}{n}} \neg v
\]


\subsubsection{Verifying SAT count for restricted BDD function}
To add an extra verification step we add a \verb|minimum_count| function that counts amount of SAT count with respect to the constraints the same as \verb|min_hi| and \verb|min_lo|. An important distinction is it only counts the paths that lead to a $\top$, it does not account for unused variables and skipped levels.
\[
\fn{minimum\_count}{n} = 
\begin{cases}
0 & n = \bot \\
1 & n = \top \\
\fn{count\_lo}{n,\lo{n}} + \fn{count\_hi}{n,\hi{n}} & otherwise\\
\end{cases}
\]
\[
\fn{count\_lo}{n,m} =
\begin{cases}
0 & \fn{min\_cost}{m} > \fn{min\_cost}{n} \\
\fn{minimum\_count}{m} & otherwise\\
\end{cases}
\]
\[
\fn{count\_hi}{n,m} = 
\begin{cases}
0 & \fn{min\_cost}{m} \ge \fn{min\_cost}{n} \\
\fn{minimum\_count}{m} & otherwise\\
\end{cases}
\]

\subsubsection{Runtime}
Each of the function that work on a BDD in its given form would visit every node once and take $O(|BDD|)$. We can reduce that by introducing caches for the BDDFunctions.
Each function that is presented contains a cache. If the given BDDFunction is already in the cache that result is returned, if not it is computed and added in the cache.
The function \verb|minimum_count|, \verb|min_config| and their respective sub functions use two caches.
The first one is the cache for the actual computation, the second cache is the \textit{cost\_cache}, that is passed to the \verb|min_cost| function.
Since the \verb|min_cost| function walks through the BDD populating the \textit{cost\_cache}, all of the sub functions in \verb|minimum_count| and \verb|min_config| use the already populated \textit{cost\_cache}.
Meaning that the those \verb|min_cost| function calls amounts to only a cache lookup.


\subsection{Results}
Per given DIMACS file we provide the calculated result in Table~\ref{tab:exercise-5-res}
\def\tabletitle #1{\multicolumn{2}{c}{\cellcolor{black}\color{white}#1}}
\newcommand\resulttable[9]{
	\begin{tabular}[t]{lp{8.5cm}}
		\tabletitle{Result for #1} \\
		(i) BDD nodes of $\llbracket \varphi \rrbracket$ & #3 \\
		(ii) Valid configurations $|V|$ & \texttt{\seqsplit{#4}} \\
		(iii) $|F\upharpoonright_{cmin}|$ & #5 \\
		(iv) Essential $|F\upharpoonright_{cmin}|$ for $S=\{x_{218},x_{311}\}$ & #6 \\
		(v) BDD nodes of $\llbracket \varphi \rrbracket_{cmin}$ & #7 \\
		(vi) Number of minimum configurations & #8 \\
		Time to run & #9 \\
\bottomrule
\end{tabular}
}
 
\begin{table}[htbp]
\caption{Results}
\centering
\begin{tabular}{c}
	\resulttable{buildroot}{565}{1000}{197929412613795104033186650529764251995382268085040059140299644922581611052953218858293468858114370125417172798513453704215828256540658252191740135040064962560}{1}{14}{567}{1}{$\sim$0.7s}
	\\
	\resulttable{busybox}{820}{3076}{341966771173357574203468511726402787658846470020778853449916660355015748144503172810981395757363854768801367747323733822923115327180062924636321384290977279599243935937986560000000000000000000000}{17}{24}{847}{243}{$\sim$1.5s}
	\\
	\resulttable{embtoolkit}{2108}{172157}{5269086459380001633643953865313843951587769385459462783768187763744059253901124797424957712252293478301274200994406811283718331120294266116775137066700763565336063060585502617122754794227972693181096959562137510764573788550603417290201944648467575812000557794560479278383157151334042379118437642922099801799434308519509535427929029197430784}{50}{111}{3441}{102114}{$\sim$2m 20s}
	\\
	\resulttable{toybox}{544}{949}{144991790900969472}{28}{51}{610}{192}{$\sim$0.6s}
	\\
	\resulttable{uClinux}{1850}{2667}{16296287810675888690147565507275025288411747149327490005089123594835050398106693649467179008}{7}{11}{1852}{1}{$\sim$5.4s}

	\label{tab:exercise-5-res}
\end{tabular}
\end{table}

\endgroup
